\chapter{Introduction}
\label{chapter:introduction}

\section{What? Why? How?}
\label{section:what}

Cardiovascular diseases are the leading cause of death worldwide, with an estimated 19.8 million deaths in 2022 \cite{who2025cvd}. Smartwatches can already screen large populations for atrial fibrillation \cite{perez2019large}, but single-lead wearable ECG is noisy and of limited diagnostic value on its own \cite{bayoumy2021smart}. Digital twins, patient-specific models updated from the patient's own data, are a promising way to personalise monitoring \cite{corralacero2020digital}. Current cardiovascular twins, however, are mostly simulation models built from clinical data \cite{coorey2022health,gillette2021framework}.

We build a lightweight twin from wearable data alone: a dynamical ECG model \cite{mcsharry2003dynamical} calibrated to the wearer's beats, combined with learned models for denoising, arrhythmia detection and heart-rate forecasting. Five agents share the twin and write a report for a clinician. The key addition is a per-beat score of how far each beat is from the patient's own normal.

\section{Paper structure and original contribution(s)}
\label{section:structure}

Chapter~\ref{section:scientificProblem} defines the problem, the five agent sub-problems and the evaluation protocol. Chapter~\ref{chapter:stateOfArt} surveys related work, and Chapter~\ref{chapter:concl} lists the next steps.
